%========================================================================================
%  SINTESI DELLA TESI (versione corretta) — Lorenzo Caivano
%  Compila questo file da solo (è indipendente dalla tesi).
%========================================================================================
\documentclass[12pt]{report}

\usepackage[utf8]{inputenc}
\usepackage[T1]{fontenc}
\usepackage[a4paper, margin=1in]{geometry}
\usepackage{setspace}
\usepackage{parskip}
\usepackage{titlesec}
\usepackage{amsmath}
\usepackage{amssymb}

% Titolo di capitolo semplice e centrato (niente "Capitolo 1")
\titleformat{\chapter}[display]
  {\normalfont\bfseries\centering}{}{0pt}{\Large}

%----------------------------------------------------------------------------------------
%  DATI DELLA TESI
%----------------------------------------------------------------------------------------
% Italiano
\newcommand{\cognomeIT}{Caivano}
\newcommand{\nomeIT}{Lorenzo}
\newcommand{\matricolaIT}{889490}
\newcommand{\titoloTesiIT}{Studio delle degenerazioni nell'analisi del rumore delle pulsar}
\newcommand{\relatoreIT}{Prof. Alberto Sesana}
\newcommand{\correlatoreIT}{Dott.ssa Irene Ferranti}
\newcommand{\dataSedutaIT}{23--25 Novembre 2026}
\newcommand{\corsoDiLaureaIT}{Fisica}
\newcommand{\recapitoTelefonicoIT}{+39 3391306307}

% Inglese
\newcommand{\surnameEN}{Caivano}
\newcommand{\nameEN}{Lorenzo}
\newcommand{\matricolaEN}{889490}
\newcommand{\thesisTitleEN}{Study of degeneracies in pulsar noise analysis}
\newcommand{\advisorEN}{Prof. Alberto Sesana}
\newcommand{\coadvisorEN}{Dr. Irene Ferranti}
\newcommand{\defenseDateEN}{23--25 November 2026}
\newcommand{\degreeCourseEN}{Physics}
\newcommand{\phoneEN}{+39 3391306307}

\begin{document}
\thispagestyle{empty}

%========================================================================================
%  SINTESI IN ITALIANO
%========================================================================================
\chapter*{Sintesi}
\begin{center}
\textbf{Nome e Cognome:} \nomeIT \, \cognomeIT \\[0.3cm]
\textbf{Matricola:} \matricolaIT \\[0.3cm]
\textbf{Titolo (trad.):} \titoloTesiIT \\[0.3cm]
\textbf{Relatore:} \relatoreIT \\[0.3cm]
\textbf{Correlatrice:} \correlatoreIT \\[0.3cm]
\textbf{Data della Prova Finale:} \dataSedutaIT \\[0.3cm]
\textbf{Corso di Laurea:} \corsoDiLaureaIT \\[0.3cm]
\textbf{Recapito Telefonico:} \recapitoTelefonicoIT \\
\end{center}

\vspace{1cm}
\sloppy

\noindent
Le \textit{millisecond pulsar} (MSP) sono stelle di neutroni in rapida rotazione che
emettono impulsi radio con una regolarità paragonabile a quella degli orologi
atomici. Un'onda gravitazionale che attraversa lo spazio tra la Terra e una pulsar
perturba il tempo di percorrenza degli impulsi, che arrivano leggermente in
anticipo o in ritardo rispetto a quanto previsto. Monitorando per decenni i
\textit{times of arrival} (TOA) degli impulsi di un insieme di pulsar, un
\textit{Pulsar Timing Array} (PTA), è quindi possibile osservare onde gravitazionali
con frequenze dell'ordine dei nanohertz ($10^{-9}$--$10^{-7}$\,Hz). \par

In questa banda la sorgente principale attesa è il fondo stocastico di onde
gravitazionali (\textit{Gravitational Wave Background}, GWB), prodotto dalla
sovrapposizione incoerente dei segnali emessi dalla popolazione di binarie di buchi
neri supermassicci. Il GWB si manifesta come un segnale a bassa frequenza comune a
tutte le pulsar e correlato tra di esse in modo caratteristico; nel 2023 diverse
collaborazioni, tra cui lo \textit{European Pulsar Timing Array} (EPTA), ne hanno
riportato una prima evidenza. \par

La rivelazione del GWB è tuttavia resa difficile dalle numerose sorgenti di rumore
presenti nei TOA degli impulsi: irregolarità nella rotazione della pulsar
(\textit{red noise}, RN), variazioni della densità di elettroni lungo la linea di
vista (\textit{dispersion measure}, DM) e diffusione degli impulsi nel mezzo
interstellare (\textit{scattering}). Questi processi producono nei residui di
timing un rumore ``rosso'', con più potenza alle basse frequenze, descritto dalla
stessa forma spettrale a legge di potenza del GWB. In una singola pulsar il
segnale gravitazionale è quindi difficilmente distinguibile dal rumore, e le
diverse componenti del modello di rumore possono risultare degeneri tra loro.
Un modello di rumore accurato è perciò essenziale, perché un GWB potrebbe essere
assorbito dal rumore della pulsar o, al contrario, imitato da esso. \par

Il mio lavoro di tesi si è concentrato sullo studio delle degenerazioni tra i
parametri del modello di rumore di una singola pulsar, PSR~J1640+2224, inclusa
nella seconda \textit{data release} dello \textit{European Pulsar Timing Array}
(EPTA DR2). Per lo studio sono stati analizzati sia i dati reali sia dati simulati
in modo realistico. \par

L'analisi è stata svolta su due versioni dei dati di J1640+2224: DR2new+, che
contiene solo le osservazioni EPTA acquisite con la strumentazione più recente,
e DR2low, che le combina con osservazioni a bassa frequenza (30--200\,MHz) dei
radiotelescopi LOFAR e NenuFAR. Il rumore bianco è stato fissato al
\textit{maximum likelihood value} ottenuto da analisi precedenti a questa tesi,
mentre i processi rossi sono descritti con sei parametri liberi, un'ampiezza e un
indice spettrale per ciascuno tra red noise, variazioni di DM e scattering.
L'inferenza bayesiana è stata condotta con Discovery, un nuovo codice basato su
JAX: il confronto tra i due dataset mostra che solo con DR2low, grazie alla
copertura in frequenza, i parametri del rumore cromatico risultano vincolati. Su
DR2low le posterior di Discovery coincidono con quelle ottenute con Enterprise, il
codice di riferimento della comunità, e mostrano una marcata degenerazione tra DM
e scattering. Questa degenerazione crea una \textit{posterior distribution}
bimodale nei parametri della DM, con un picco che contiene il \textit{maximum
likelihood value} ad ampiezza bassa ($\log_{10}A \sim -15$) e \textit{spectral
index} $\gamma \sim 4$, e un picco ad ampiezza più elevata e \textit{spectral
index} più basso. Inoltre genera una \textit{joint posterior} fra DM e scattering
con forma a ``L'', sintomo, appunto, del fatto che le due componenti non vengono
correttamente separate. Attraverso l'analisi di simulazioni realistiche,
l'obiettivo di questa tesi è capire se componenti di rumore non modellate, ma
presenti nei dati, possano generare questa degenerazione. \par

Assumendo che un GWB sia presente nei dati, questo rappresenta un red noise
acromatico aggiuntivo rispetto a quello generato dalle irregolarità nella
rotazione della pulsar. Per studiarne l'effetto ho simulato dati con le stesse
caratteristiche di J1640+2224, iniettando un fondo con $\gamma = 13/3$ e ampiezze
comprese tra $\log_{10}A = -15.5$ e $-14$ (compatibili con le attuali misure dei
PTA), e li ho analizzati con il modello a sei parametri (RN, DM e scattering)
privo del termine di GWB. Per $\log_{10}A = -14$ il segnale provoca una
degenerazione fra DM e scattering con caratteristiche simili a quelle osservate
nei dati, seppur con significatività molto inferiore. \par

Ho inoltre incluso un termine di \textit{band noise}, un processo rosso che agisce
solo sui TOA della banda LOFAR. Analizzando simulazioni in cui il band noise è
iniettato, ma non modellato, ottengo che questo rumore genera un bias in tutti i
processi, ma soprattutto nella \textit{posterior distribution} della DM, che si
sposta verso ampiezze più alte e \textit{spectral index} più bassi rispetto ai
valori iniettati. La presenza di un band noise non modellato crea dunque un picco
ad ampiezze e \textit{spectral index} compatibili con il picco secondario delle
\textit{posterior distribution} ottenute con i dati reali. Inoltre, le simulazioni
mostrano che, a causa dell'elevato rumore bianco dei dati LOFAR, questo processo è
rilevabile solo per ampiezze $\log_{10}A \gtrsim -12.5$; sui dati reali se ne
ottiene solo un limite superiore, simile a quello che si ottiene sui dati simulati
iniettando un band noise con $\log_{10}A \sim -12.5$. Di conseguenza, la presenza
di band noise nei dati reali non può essere esclusa e la sua ricerca sarà oggetto
di studi futuri. \par

In conclusione, l'analisi di J1640+2224 mostra che Discovery riproduce i
risultati di Enterprise e che la stima del rumore di questa pulsar è dominata
dalla degenerazione tra le due componenti cromatiche, DM e scattering, che solo
una buona copertura in frequenza radio permette di ridurre. Le simulazioni
indicano inoltre che un GWB di ampiezza paragonabile a quella oggi osservata, se
non modellato, può causare una degenerazione fra DM e scattering, e che un band
noise non modellato può generare un forte bias nei parametri della DM. Questi
risultati sottolineano l'importanza di una modellizzazione dettagliata del rumore
delle singole pulsar, presupposto essenziale per la rivelazione del GWB con i PTA.
\par

%========================================================================================
%  ABSTRACT IN INGLESE
%========================================================================================
\newpage
\chapter*{Abstract}
\begin{center}
\textbf{Name and Surname:} \nameEN \, \surnameEN \\[0.3cm]
\textbf{Student ID:} \matricolaEN \\[0.3cm]
\textbf{Title:} \thesisTitleEN \\[0.3cm]
\textbf{Advisor:} \advisorEN \\[0.3cm]
\textbf{Co-advisor:} \coadvisorEN \\[0.3cm]
\textbf{Date of Dissertation:} \defenseDateEN \\[0.3cm]
\textbf{Degree Course:} \degreeCourseEN \\[0.3cm]
\textbf{Phone Number:} \phoneEN \\
\vspace{1cm}
\end{center}
\sloppy

\noindent
Millisecond pulsars (MSPs) are rapidly rotating neutron stars that emit radio
pulses with a regularity comparable to that of atomic clocks. A gravitational
wave crossing the space between the Earth and a pulsar perturbs the travel time
of the pulses, which arrive slightly earlier or later than predicted. By
monitoring for decades the times of arrival (TOAs) of the pulses from a set of
pulsars, a Pulsar Timing Array (PTA), it is therefore possible to observe
gravitational waves with frequencies of the order of nanohertz
($10^{-9}$--$10^{-7}$\,Hz). \par

In this band, the main expected source is the stochastic gravitational wave
background (GWB), produced by the incoherent superposition of the signals
emitted by the population of supermassive black hole binaries. The GWB appears
as a low-frequency signal common to all pulsars and correlated among them in a
characteristic way; in 2023 several collaborations, including the European
Pulsar Timing Array (EPTA), reported first evidence for it. \par

The detection of the GWB is, however, made difficult by the many noise sources
affecting the pulse TOAs: irregularities in the rotation of the pulsar (red
noise, RN), variations of the electron density along the line of sight
(dispersion measure, DM), and scattering of the pulses in the interstellar
medium. These processes produce a ``red'' noise in the timing residuals, with
more power at low frequencies, described by the same power-law spectral shape as
the GWB. In a single pulsar the gravitational signal is therefore hard to
distinguish from the noise, and the different components of the noise model can
be degenerate with one another. An accurate noise model is thus essential, since
a GWB could be absorbed by the pulsar noise or, conversely, mimicked by it. \par

My thesis work focused on the study of the degeneracies among the parameters of
the noise model of a single pulsar, PSR~J1640+2224, included in the second data
release of the European Pulsar Timing Array (EPTA DR2). Both the real data and
realistically simulated data were analysed. \par

The analysis was carried out on two versions of the J1640+2224 data: DR2new+,
which contains only the EPTA observations taken with the most recent
instrumentation, and DR2low, which combines them with low-frequency observations
(30--200\,MHz) from the LOFAR and NenuFAR radio telescopes. The white noise was
fixed to the maximum likelihood values obtained in analyses prior to this
thesis, while the red processes were described by six free parameters, an
amplitude and a spectral index for each of red noise, DM variations and
scattering. The Bayesian inference was performed with Discovery, a new code
based on JAX: the comparison between the two datasets shows that only with
DR2low, thanks to its frequency coverage, the parameters of the chromatic noise
are constrained. On DR2low, the posteriors obtained with Discovery coincide with
those obtained with Enterprise, the reference code of the community, and show a
strong degeneracy between DM and scattering. This degeneracy produces a bimodal
posterior distribution for the DM parameters, with one peak, containing the
maximum likelihood value, at low amplitude ($\log_{10}A \sim -15$) and spectral
index $\gamma \sim 4$, and a second peak at higher amplitude and lower spectral
index. It also produces an ``L''-shaped joint posterior between DM and
scattering, a sign that the two components are not correctly separated. Through
the analysis of realistic simulations, the aim of this thesis is to understand
whether noise components that are present in the data but not modelled can
generate this degeneracy. \par

If a GWB is present in the data, it represents an additional achromatic red
noise on top of the one generated by the rotational irregularities of the
pulsar. To study its effect, I simulated data with the same characteristics as
J1640+2224, injecting a background with $\gamma = 13/3$ and amplitudes between
$\log_{10}A = -15.5$ and $-14$ (consistent with current PTA measurements), and
analysed them with the six-parameter model (RN, DM and scattering) without a
GWB term. For $\log_{10}A = -14$, the signal induces a degeneracy between DM and
scattering with features similar to those observed in the data, albeit with a
much lower significance. \par

I also included a band noise term, a red process acting only on the TOAs of the
LOFAR band. By analysing simulations in which the band noise is injected but not
modelled, I find that this noise biases all the processes, and most of all the
DM posterior distribution, which shifts towards higher amplitudes and lower
spectral indices with respect to the injected values. An unmodelled band noise
therefore produces a peak at amplitudes and spectral indices consistent with the
secondary peak of the posterior distributions obtained from the real data.
Moreover, the simulations show that, because of the high white noise of the
LOFAR data, this process is detectable only for amplitudes
$\log_{10}A \gtrsim -12.5$; on the real data only an upper limit is obtained,
similar to the one obtained on simulated data with an injected band noise of
$\log_{10}A \sim -12.5$. Consequently, the presence of band noise in the real
data cannot be excluded, and its search will be the subject of future work. \par

In conclusion, the analysis of J1640+2224 shows that Discovery reproduces the
results of Enterprise and that the noise estimate for this pulsar is dominated
by the degeneracy between the two chromatic components, DM and scattering, which
only a good radio-frequency coverage can reduce. The simulations further
indicate that a GWB with an amplitude comparable to the one currently observed
can, if not modelled, cause a degeneracy between DM and scattering, and that an
unmodelled band noise can strongly bias the DM parameters. These results
highlight the importance of a detailed modelling of the noise of individual
pulsars, an essential requirement for the detection of the GWB with PTAs. \par

\end{document}
